\section{Discussion}\label{sec:discussion}

\subsection{Architectural and RI Implications}
The main architectural lesson is that the gateway should be treated as a semantic boundary, not as a transparent relay. Constrained devices can remain focused on protected communication, local execution, and evidence emission, while the gateway handles identity validation, protocol translation, and status propagation toward Kubernetes. This separation avoids pushing cloud-facing semantics into firmware and gives the control plane a stable interface even if the southbound protocol evolves.

The evaluation confirms that protocol success alone is not sufficient: enrollment, liveness, and deployment evidence must be interpreted consistently before they become trustworthy control-plane state. This supports the central design decision to separate desired state, gateway-derived evidence, and controller reconciliation. Future extensions such as multi-gateway placement, staged rollout, or retry policies should preserve this separation rather than allowing multiple components to write conflicting interpretations of device status.

For sustainable digital research infrastructures, the same separation creates explicit attachment points for energy and carbon policies. Placement decisions can be expressed as desired state, gateway and device observations can be treated as evidence, and sustainability controllers can consume the same status model without requiring constrained devices to run additional cloud-native agents. The measured baselines in Section~\ref{sec:sustainability-metrics}, byte-scale southbound traffic, single-digit-MiB gateway mediation, and no kubelet on the edge, suggest that subsequent green-experimentation studies can focus on policy and instrumentation rather than on redesigning the orchestration substrate.

%{\color{red}For the scope of sustainable digital research infrastructures, the main outcome is therefore methodological rather than directly energetic: the paper identifies where orchestration evidence, device liveness, gateway mediation, and resource overhead can be captured consistently before energy-aware or carbon-aware decisions are automated.}

\begingroup
\color{red}
The practical implication for research infrastructures is that orchestration evidence becomes a reusable experimental object. A placement policy, for example, can be evaluated by inspecting which Application resources were scheduled to which devices, which gateway mediated the execution, whether heartbeat evidence remained fresh, and whether the reported execution outcome matched the intended state. The same record can later be joined with external measurements such as site power, carbon-intensity annotations, thermal constraints, or gateway utilization. The current prototype does not optimize those signals, but it provides the control-plane hooks that a sustainability-aware controller would need in order to make such optimization auditable.

%Removed (ridondanza): The gateway boundary also simplifies reproducibility. A reproduced experiment does not need the embedded device to expose Kubernetes-native behavior; it only needs the same southbound protocol semantics, the same CRD intent, and the same interpretation of evidence. This is useful for multi-site RIs, where physical boards, local networks, and energy meters may differ across laboratories. The orchestration claim can be retested with emulation, while hardware-specific claims can be added as separate measurement layers.

\endgroup

\subsection{Validity and Scope}
The validation targets system-level orchestration behavior rather than real-time performance, radio behavior, or sustainability optimization. Emulation and scripted trials are appropriate for exercising enrollment, deployment, status convergence, and reconciliation semantics, while physical boards, wireless links, and power meters require additional characterization. The results should therefore be interpreted as evidence that the orchestration model is coherent under the evaluated workflow.

% Removed: External validity is bounded by the single-node K3s deployment, one gateway path, and a representative minimal workload.
{\color{red}External validity is bounded by the single-node K3s deployment, one gateway path, and a minimal WASM workload selected to isolate orchestration overhead rather than application-computation cost.} Larger deployments introduce gateway contention, distributed failure modes, network partitions, and device mobility. % Removed: A supplementary scaling attempt with three Renode machines on the shared TAP test bed kept only one TLS session active in the gateway runtime view; fleet-scale energy characterization therefore requires isolated network segments or physical boards.
%Removed (compressione): {\color{red}Concurrency is exercised directly rather than assumed: three emulated devices attached simultaneously hold three distinct TLS sessions in the gateway runtime view and each receive their own deployment (Section~\ref{sec:concurrent-devices}). The 50-board registry exercise remains a registry-state scaling check, since those boards are registry entries rather than TLS-connected devices, and the concurrency evidence stops at the three devices actually measured on a single host: it does not license extrapolation to hundreds of devices, to multiple gateways, or to fleet-scale energy characterization, all of which remain open.}
{\color{red}Concurrency is exercised rather than assumed: three emulated devices hold
three distinct TLS sessions and each receives its own deployment
(Section~\ref{sec:concurrent-devices}). That evidence stops at three devices on one host,
and the 50-board exercise scales registry state rather than sessions, so neither extends to
multiple gateways or to fleet-scale energy characterization.} Construct validity is also scoped: the study measures %Removed (M11): "secure attachment"
\revD{identity-bound attachment}, observable liveness, deployment completion, latency, wire sizes, firmware footprint, and control-plane RAM/CPU samples, while direct, hardware-validated energy consumption and long-duration reliability remain outside the current evaluation%Removed (compressione): {\color{blue}; Section~\ref{sec:energy-instrumentation} reports a computational-energy instrumentation pass on a representative WASM workload using exact cgroup CPU-time as the primary signal and a model-based Kepler host-power estimate as a secondary one, as a step toward that validation.}
{\color{blue}; Section~\ref{sec:energy-instrumentation} is a step toward the latter, with exact cgroup CPU-time as the primary signal.}

\begingroup
\color{blue}
\paragraph{Scope of the energy signal.} cgroup CPU-time (\texttt{cpu.stat}) is the only quantity in Section~\ref{sec:energy-instrumentation} that is an exact, reproducible measurement rather than a model output, and it alone should be treated as load-bearing for any energy claim in this paper. 
%The Kepler wattage is a secondary, corroborating signal: this VM exposes no RAPL counters, so every reported watt is the output of Kepler's regression-based power model rather than a hardware reading, independent of how large the sample is. The Joule and carbon figures in Table~\ref{tab:workload-cpu-time} are a further derived layer on top of the measured CPU-time, using a literature power coefficient, and are reported as illustrative estimates rather than measurements. 
A hardware-validated absolute wattage for this workload, on a host that combines a working k3s control plane with exposed RAPL counters, remains future work.
\endgroup

Conclusion validity is strengthened by the transactional success criterion: enrollment, heartbeat visibility, and deployment must succeed within the same trial. This avoids overstating correctness from isolated protocol stages. Broader claims about optimality, scalability, or comparative advantage require additional experiments against alternative orchestration designs and physical-device deployments.

\begingroup
\color{red}
\subsection{Scalability and Failure Modes}
The main scalability pressure point is the gateway, because it concentrates southbound sessions, protocol parsing, and status propagation. This is an intentional trade-off: moving the boundary into the gateway avoids embedding Kubernetes behavior in every device, but it also means that future deployments must treat gateways as schedulable and observable resources. Multi-gateway operation should therefore add explicit placement policies, per-gateway capacity metadata, and failover rules rather than relying on implicit network reachability. In such a design, a Device resource would declare or discover a preferred gateway, while Application reconciliation would resolve the current gateway association before attempting deployment.

%Removed (compressione): The evaluated system also identifies four categories of faults. Transport faults include TLS interruption, device reboot, and heartbeat loss. Protocol faults include malformed CBOR records, identity mismatch, stale session state, and deployment results that cannot be associated with the requested Application. Control-plane faults include delayed status patches, controller restarts, and conflicting updates from independent reconcilers. Application faults include WASM load failure, runtime error, or an execution result that does not satisfy the expected policy. Treating these cases separately is important because each one requires a different recovery action: reconnect, reject enrollment, retry deployment, preserve desired state, or surface an explicit error to the operator.
\revD{Four fault categories require distinct recovery actions and are therefore kept
separate: transport faults (TLS interruption, reboot, heartbeat loss) call for reconnection;
protocol faults (malformed CBOR, identity mismatch, stale sessions, unassociable deployment
results) for rejection; control-plane faults (delayed patches, controller restarts,
conflicting reconcilers) for preserving desired state; and application faults (WASM load
failure, runtime error, policy violation) for surfacing an explicit error.}

%Removed (artefatto): The above-stated considerations also limits the interpretation of the present results.
\revD{These considerations also limit the interpretation of the present results.} The 100-trial campaign shows that the nominal enrollment-to-deployment transaction is coherent under the evaluated emulated conditions. It does not establish maximum gateway capacity, behavior under network partitions, real-time scheduling guarantees, or recovery time after simultaneous failures. Those are natural next experiments once the single-device protocol and reconciliation invariants are fixed.
\endgroup

\subsection{Reproducibility}
The evaluation is designed to be repeatable on a commodity Linux host with Docker and K3s. Renode makes the embedded side deterministic enough for regression-style validation, although it remains complementary to future physical-board experiments. Evidence traceability is maintained across multiple layers: gateway logs, API-visible resource states, and Kubernetes CRD status describe the same workflow from different perspectives, so that protocol success and control-plane convergence are not inferred from a single eventually consistent field.

\begingroup
\color{red}
For a replication package, the most important artifacts are therefore not only the source code and container images, but also the CRD manifests, firmware configuration, WASM module, trial harness, raw latency records, and the scripts that generate the figures. Reporting these artifacts together would allow another researcher to distinguish implementation regressions from environmental variation. A failed replication could then be localized to a specific layer: Kubernetes resource creation, gateway mediation, TLS/CBOR communication, embedded execution, or post-processing of the measurement records. {\color{red}If institutional constraints prevent publishing all raw records, the artifact statement should separate public assets, such as source code, manifests, and scripts, from records available on request.}
\endgroup
